% !TeX encoding = UTF-8
% !TeX spellcheck = en_GB
% !TeX root = ../Thesis_AndreGuerra.tex
%%%%%%%%%%%%%%%%%%% Thesis_AG_Preamble.tex %%%%%%%%%%%%%%%%%%
%
% PhD thesis - Preamble
% André Guerra
%
% 04/Apr/2017
%
%%%%%%%%%%%%%%%%%%%%%%%%%%%%%%%%%%%%%%%%%%%%%%%%%%%%%%%%%%%%%

%------------------------------------------------------------
%% General ------------------------------------------------%%

% Comments and review packages
\usepackage{lipsum}     % For draft text
\usepackage[linecolor=red,figcolor=white,colorinlistoftodos,%
]{todonotes}                                                % TEMP command to review
\newcommand{\review}[1]{{\textcolor{red}{\framebox{#1}}}}	% TEMP command to review

\pdfminorversion = 7

%%% Book configurations for report class
%%\makeatletter % changes the catcode of @ to 11
    %%\newcommand\frontmatter{%
        %%\cleardoublepage
        %%%\@mainmatterfalse
        %%\pagenumbering{roman}
    %%}
%%
    %%\newcommand\mainmatter{%
        %%\cleardoublepage
        %%% \@mainmattertrue
        %%\pagenumbering{arabic}
    %%}
%%
    %%\newcommand\backmatter{%
        %%\if@openright
        %%\cleardoublepage
        %%\else
        %%\clearpage
        %%\fi
        %%% \@mainmatterfalse
    %%}
%%\makeatother % changes the catcode of @ back to 12

%------------------------------------------------------------
%% Language -----------------------------------------------%%
\usepackage[main=UKenglish,portuguese]{babel}   % Language hyphenation and typographical rules
\usepackage[T1]{fontenc}                        % Use 8-bit encoding that has 256 glyphs
\usepackage[utf8]{inputenc}                     % Accept different input encodings.

\babeltags{pt = portuguese}
%\textpt{Portuguese text}
%\begin{pt}
%    Portuguese text
%\end{pt}

%------------------------------------------------------------
%% Page geometry and others & TOC -------------------------%%
\usepackage{geometry}	
\geometry{verbose,
    a4paper,
    twoside,                        % switches on twoside mode with left and right margins swapped on verso pages
    includehead,                    % includes the head of the page
    includefoot,                    % includes the foot of the page
    textheight      = 674pt,        % the height of body (the text area)
    textwidth       = 445pt,        % the width of body (the text area)
    headheight      = 40.0pt,       % height of header
    headsep         = 21.68121pt,   % separation between header and text
    footskip        = 27.46295pt,   % distance separation between baseline of last line of text and baseline of footer
    marginparsep    = 7.0pt,        % separation between body and marginal notes
    marginparwidth  = 70pt,         % width of the marginal notes
    inner           = 22mm          % left margin (for oneside) or inner margin (for twoside)
    %
    %outer=25mm,
    %top=30mm,
    %bottom=20mm,
    %headheight=15pt
}
\marginparpush = 5.0pt

% Depth of TOC
\setcounter{secnumdepth}{3}
\setcounter{tocdepth}{2}

% Control over table of contents, figures, etc
\usepackage{tocloft}

% General purpose hierarchical index generator
\usepackage{makeidx}
\makeindex

% Nomenclature package
\usepackage[english]{nomencl}
\makenomenclature
\usepackage{ifthen}
\renewcommand{\nomgroup}[1]{%
    \ifthenelse{\equal{#1}{R}}{\item[\textbf{Roman Symbols}]}{%
        \ifthenelse{\equal{#1}{G}}{\item[\textbf{Greek Symbols}]}{%
            \ifthenelse{\equal{#1}{A}}{\item[\textbf{Acronyms/Abbreviations}]}{%
                \ifthenelse{\equal{#1}{T}}{\item[\textbf{Superscripts}]}{%
                    \ifthenelse{\equal{#1}{S}}{\item[\textbf{Subscripts}]}{%
                        \ifthenelse{\equal{#1}{X}}{\item[\textbf{Other Symbols}]}{%
                            \ifthenelse{\equal{#1}{E}}{\item[\textbf{Electronic Symbols}]}
                            {}
                        }% matches Electronic Roman symbols > E
                    }% matches mathematical symbols > X
                }% matches Subscripts           > T
            }% matches Superscripts         > R
        }% matches Abbreviations        > A
    }% matches Greek Symbols        > G
}% matches Roman Symbols        > R
%% To add some text before printing the list of symbols
\renewcommand{\nompreamble}{This is a list of several symbols and abreviations used within the body of the document.}
%
% This will add the units to the symbols
%----------------------------------------------
\newcommand{\nomunit}[1]{\renewcommand{\nomentryend}{\hspace*{\fill}#1}}
%----------------------------------------------

% Appendix
\usepackage[page,titletoc]{appendix}
\renewcommand{\appendixpagename}{Appendix}

% Include PDF documents in LaTeX
\usepackage{pdfpages}

%------------------------------------------------------------
%% Font & Paragraph ---------------------------------------%%
\usepackage{sansmathfonts}                      % For math fonts
\usepackage[scaled]{helvet}                     % Font family
\renewcommand{\familydefault}{\sfdefault}
\renewcommand{\rmdefault}{\sfdefault}

\usepackage{sectsty}    % change the style of any or all of LATEX's sectional headers
\allsectionsfont{\sfdefault}

\usepackage{fncychap}   % for changing chapter level headings
% CHAPTER font size
\ChNameVar{\fontsize{14}{16}\usefont{OT1}{phv}{m}{n}\selectfont}
\ChNumVar{\fontsize{60}{62}\usefont{OT1}{ptm}{m}{n}\selectfont}
\ChTitleVar{\Huge\bfseries\rm} \ChRuleWidth{1pt}
% Font size of number X
\ChNumVar{\raggedleft\fontsize{20pt}{25pt}\usefont{OT1}{ptm}{m}{n}\selectfont}
\ChTitleVar{\raggedright\rm\fontsize{20pt}{25pt}\bfseries\selectfont}
\ChNameAsIs
\ChNameVar{[1]}

% Chapter style
\ChNameUpperCase
\ChNameVar{\Huge\sf\bf\flushleft}
\ChNumVar{\Huge\sf\bf}
\ChTitleVar{\huge\bf\raggedright}
% Section Style
\sectionfont{\Large\raggedright}
\subsectionfont{\large\raggedright}
\subsubsectionfont{\normalsize\raggedright}

% settings for paragraphs
\setlength{\parskip}{2mm}
% defining the space between lines (1.5 is for double-space)
\usepackage{setspace}
\onehalfspacing
%%\renewcommand{\baselinestretch}{1.5}

% AMS mathematical facilities for LaTeX.
\usepackage{amsmath,amsfonts,amssymb,bm}

\usepackage{eurosym}            % Euro sign

\usepackage[per-mode = symbol,binary-units = true]{siunitx}
\sisetup{group-minimum-digits = 4}  % Group with a minimum of 4 numbers (i.e. 1 000)
\DeclareSIUnit{\EUR}{\text{\euro}}
\DeclareSIUnit[number-unit-product = {}]{\degree}{\SIUnitSymbolDegree}

\usepackage{cancel}     % To have symbols crossed on equations
\usepackage{gensymb}    % Provides generic commands (\degree, \celsius, \perthousand, \micro and \ohm)
\usepackage{xfrac}      % allows the user to typeset fractions in the form 'n/d' (\sfrac{1}{2})

% Prevent hyphenation and paragraph breaks
%\hyphenpenalty=10000
%\brokenpenalty=5000
%\clubpenalty=5000
%\widowpenalty=5000

\flushbottom    % makes all text pages the same height

\usepackage[activate={true,nocompatibility},final,tracking=true,kerning=true,spacing=true,factor=1100,stretch=10,shrink=10]{microtype}
% activate={true,nocompatibility} - activate protrusion and expansion
% final - enable microtype; use "draft" to disable
% tracking=true, kerning=true, spacing=true - activate these techniques
% factor=1100 - add 10% to the protrusion amount (default is 1000)
% stretch=10, shrink=10 - reduce stretchability/shrinkability (default is 20/20)
\microtypecontext{spacing=nonfrench}

\SetExtraKerning[unit=space]
{encoding={*}, family={*}, series={*}, size={footnotesize,small,normalsize}}
{\textendash={400,400}, % en-dash, add more space around it
    "28={ ,150}, % left bracket, add space from right
    "29={150, }, % right bracket, add space from left
    \textquotedblleft={ ,150}, % left quotation mark, space from right
    \textquotedblright={150, }} % right quotation mark, space from left

\SetExtraKerning[unit=space]
{encoding={*}, family={qhv}, series={b}, size={large,Large}}
{1={-200,-200}, 
    \textendash={400,400}}
%%
%%\SetTracking{encoding={*}, shape=sc}{40}

\usepackage{enumitem}   % Customized lists
\setlist{noitemsep}     % Make itemize lists more compact

% Numbering of equations, figures and tables
\numberwithin{equation}{chapter}
\numberwithin{figure}{chapter}
\numberwithin{table}{chapter}

\usepackage[explicit]{titlesec}
\titleformat{name=\paragraph,numberless}[runin]{\normalfont\normalsize\bfseries}{}{15pt}{#1.}

%------------------------------------------------------------
%% Graphics & Figures -------------------------------------%%
\usepackage{xcolor}         % Package to define colours
\definecolor{FCUPBlue}{RGB}{147,191,235}    % FCUP color

\usepackage{graphicx}   % Enhanced support for graphics.

% Custom captions under/above floats in tables or figures
\usepackage{caption}
\usepackage{subcaption}
\renewcommand{\figurename}{Fig.} %to support older versions of captions.sty
\captionsetup{format=plain,indention=0.5cm,labelsep=endash,justification=centerlast,font=footnotesize,labelfont=bf,textfont=it,tableposition=top}

\usepackage[section]{placeins}  % ensure all floats for a section appear before the next \section command
\usepackage{flafter}            % force floats to appear after they are defined

\usepackage{epstopdf}   % To include eps figures

\usepackage{booktabs}   % For professional looking tables
\usepackage{multirow}   % Create tabular cells spanning multiple rows
\usepackage{multicol}   % Create tabular cells spanning multiple columns
\usepackage{rotating}   % To rotate tables and cell text
\usepackage{threeparttable} % tables with footnotes

% To use pgf-tikz figures
\usepackage{pgf,tikz}
\usepackage{mathrsfs}
\usetikzlibrary{arrows}
\usepackage{pgfplots}
\pgfplotsset{compat=1.15}

% Path to the figures directory
\graphicspath{
    {1_Figs/}
    {01_FrontMatter/Figs/}
    {10_Introduction/Figs/}
    {20_SmallSatellites/Figs/}
    {30_Oceanography/Figs/}
    {40_IntegratedScenario/Figs/}
    {50_ThermalAnalysis/Figs/}
    {80_Appendices/Figs/}}

%------------------------------------------------------------
%% Links and Bookmarks ------------------------------------%%
\usepackage[hyphens]{url}   % To include urls on document and bibliography

% 'hyperref' package
\usepackage{hyperref} % enhance documents that are to be output as HTML and PDF
\hypersetup{
    colorlinks          = true,         % color text of links and anchors, eliminates borders around links
    allcolors           = black,        % color for anchor text (black)
    bookmarksnumbered   = true,         % put section numbers in bookmarks
    bookmarksopen       = true,         % don't expand bookmarks
    breaklinks          = true,         % allow links to break over lines 
    %citecolor           = FCUPblue,     % color of citation links
    %backref             = true,         % do bibliographical back references
    %
    pdfdisplaydoctitle  = true,         % display document title instead of file name in title bar
    pdfpagelayout       = OneColumn,    % set layout of PDF pages
    pdfstartview        = FitH,         % starting view of PDF document
    pdfpagemode         = UseNone,      % set default mode of PDF display (UseOutlines)
    %
    pdftitle            = {PhD Thesis Andre Guerra},
    pdfauthor           = {Andre Guerra},
    pdfsubject          = {PhD thesis of Andre Guerra.}%,
    %pdfkeywords         = {\@keywords}    
}

\usepackage[figure,table]{hypcap}       % Adjusting the anchors of captions.

\usepackage[open,openlevel=2]{bookmark} % For advanced bookmark options

%------------------------------------------------------------
%% Bibliography -------------------------------------------%%
\usepackage[square,comma,numbers,sort&compress,longnamesfirst]{natbib}

\usepackage{notoccite}  % Prevent trouble from citations in table of contents, etc

\usepackage{doi}        % hyperlink DOI numbers to dx.doi.org

%------------------------------------------------------------
%% Headers and Footers ------------------------------------%%
\usepackage{fancyhdr}

\pagestyle{fancy}   % Defines the page style to fancy
\fancyhf{}          % Clear headers and footers

%%%% Headers
% Remove headers and footers on blank pages
\makeatletter % changes the catcode of @ to 11
\def\cleardoublepage{\clearpage\if@twoside \ifodd\c@page\else
    \hbox{}
    %\vspace*{\fill}
    %\begin{center}
    %This page intentionally contains only this sentence.
    %\end{center}
    %\vspace{\fill}
    \thispagestyle{empty}
    \newpage
    \if@twocolumn\hbox{}\newpage\fi\fi\fi}
\makeatother % changes the catcode of @ back to 12

% Define chapter name in smallcaps at header
\renewcommand{\chaptermark}[1]{\markboth{#1}{a}}     
\renewcommand{\sectionmark}[1]{\markright{\thesection\ #1}}

% Remove header and footer lines
\renewcommand{\headrulewidth}{0pt}
\renewcommand{\footrulewidth}{0pt}

%%%% Change chapter top margin
\makeatletter % changes the catcode of @ to 11
% For the case \chapter
\renewcommand*{\@makechapterhead}[1]{%
    \vspace*{1\p@}%
    {\parindent \z@ \raggedright \normalfont
        \ifnum \c@secnumdepth >\m@ne
        \if@mainmatter    % Fix for frontmatter, mainmatter, and backmatter 040920
        \DOCH
        \fi
        \fi
        \interlinepenalty\@M
        \if@mainmatter  % Fix for frontmatter, mainmatter, and backmatter 060424
        \DOTI{#1}%
        \else%
        \DOTIS{#1}%
        \fi
}}
% For the case \chapter*
\renewcommand*{\@makeschapterhead}[1]{%
    \vspace*{1\p@}%
    {\parindent \z@ \raggedright
        \normalfont
        \interlinepenalty\@M
        \DOTIS{#1}
        \vskip 10\p@
}}
\makeatother % changes the catcode of @ back to 12

\fancypagestyle{myfancy}{
    % Odd pages with right side header
    \fancyhead[RO]{
        \begin{tabular}{r|c} %%% REMOVE (\textbf{DRAFT \today}) for final version!!!
            {\fontsize{8pt}{8pt}\selectfont FCUP (\textbf{DRAFT \today})}                   & \multirow{2}{*}{\fontsize{8pt}{10pt}\selectfont\thepage} \\
            {\fontsize{8pt}{8pt}\selectfont\nouppercase\rightmark}  & \\
        \end{tabular}%
    }
    %
    % Even pages with left side header
    \fancyhead[LE]{
        \begin{tabular}{c|l} %%% REMOVE (\textbf{DRAFT \today}) for final version!!!
            \multirow{2}{*}{\fontsize{8pt}{10pt}\selectfont \thepage}           & {\fontsize{8pt}{8pt}\selectfont FCUP (\textbf{DRAFT \today})} \\
            & {\fontsize{8pt}{8pt}\selectfont\nouppercase\leftmark} \\
        \end{tabular}%
    }
    %
    % Odd (Even) pages left (right) side header
    \fancyhead[LO,RE]{
        \raisebox{-0.4\height}{\includegraphics[height=30pt, keepaspectratio=true]{FCUP_eps}}
    }
}

\fancypagestyle{plain}{%
    \fancyhead[RO,LE]{{\fontsize{8pt}{10pt}\selectfont\thepage}}
    \fancyhead[LO,RE,CE,CO]{}
}
\fancypagestyle{myfrontmatter}{%
    \fancyhead[RO,LE]{{\fontsize{8pt}{10pt}\selectfont\thepage}}
    \fancyhead[LO,RE,CE,CO]{}
}

% Footnotes
\usepackage{footnote}
% Footnote style
\renewcommand{\thefootnote}{\roman{footnote}}
\interfootnotelinepenalty=10000

%------------------------------------------------------------
%% Customisations -----------------------------------------%%

% Change the way nomenclature is presented
\makeatletter % changes the catcode of @ to 11
\@ifundefined{chapter}
{\def\wilh@nomsection{section}}
{\def\wilh@nomsection{chapter}}

\def\thenomenclature{%
    \begin{multicols*}{2}[%
        \csname\wilh@nomsection\endcsname*{\nomname}
        \if@intoc\addcontentsline{toc}{\wilh@nomsection}{\nomname}\fi
        \nompreamble]
        \setlength{\columnsep}{1cm}
        \setlength{\columnseprule}{1pt}
        \list{}{%
            \labelwidth\nom@tempdim
            \leftmargin\labelwidth
            \advance\leftmargin\labelsep
            \itemsep\nomitemsep
            \let\makelabel\nomlabel
        }%
    }
    
    \def\endthenomenclature{%
        \endlist
    \end{multicols*}
    \nompostamble}
\makeatother % changes the catcode of @ back to 12

%------------------------------------------------------------
%% Standarisation -----------------------------------------%%

% Hyperref package names for referecing parts of text
% Macro                 Default
%\figurename            Figure          % Reference a figure
%\tablename             Table           % Reference a table
%\appendixname          Appendix
%\equationname          Equation
%\chaptername           chapter
%\sectionname           section         % Reference a section
%\subsectionname        subsection
%\subsubsectionname     subsubsection
%\paragraphname         paragraph
%\Hfootnotename         footnote

\addto\extrasUKenglish{%
    \renewcommand{\chapterautorefname}{Chapter}%
}
\addto\extrasUKenglish{%
    \renewcommand{\sectionautorefname}{Section}%
}
\addto\extrasUKenglish{%
\renewcommand{\subsectionautorefname}{Section}%
}
\addto\extrasUKenglish{%
    \def\equationautorefname~#1\null{Equation~(#1)\null}%
}


% end of file Thesis_AndreGuerra.tex
%%%%%%%%%%%%%%%%%%%%%%%%%%%%%%%%%%%%%%%%%%%%%%%%%%%%%%%%%%%%%%%%%%%